% !TeX encoding = UTF-8
% !TeX spellcheck = en_GB
% !TeX root = ../thesis.tex


\chapter{Introduction}
\label{chapter:introduction}

In order to start using this template, change the values inside \texttt{thesis.tex} to your liking. \emph{It is strongly advised to only change what is necessary (language, names etc.).} Please also remember to remove the "REMOVE ME LATER" part when you are done. Below you can find some examples of important \LaTeX-commands.

This is a cited text~\cite{taubmann2016CloudPhylactor}.

\todo{Add more content}

This is a reference to \Cref{chapter:conclusions,sec:some-section}.

We can also add code such as in \Cref{lst:log}.

\begin{code}
    \captionof{listing}{Captions for listings are usually placed above. Java code calculating the logarithm of a given number with respect to a given base.}
    \label{lst:log}
    \begin{minted}{java}
public static double log(double x, double base) {
    return Math.log(x) / Math.log(base);
}
    \end{minted}
\end{code}

You can also include and reference figures and tables such as \Cref{fig:test} and \Cref{tab:test}. Notice how the table cannot be placed using the "h" specifier and thus uses the "t" specifier instead?

\begin{figure}[ht!]
    \centering
    \includegraphics[width=0.5\textwidth]{img/logouni.png}
    \caption{Captions for figures are usually placed below. The German logo of the University of Passau.}
    \label{fig:test}
\end{figure}

URLs can be added for example like this: \url{https://www.fim.uni-passau.de/technische-informatik/}.

\begin{table}[ht!]
    \centering
    \caption{Captions for tables are usually placed above. $\phi$ denotes the Euler totient function, by the way.}
    \label{tab:test}
    \begin{tabular}{c|c}
        $x$ & $\phi(x)$ \\
        \hhline{=|=}
        $1$ & $1$ \\
        \hline
        $2$ & $1$ \\
        \hline
        $3$ & $2$ \\
        \hline
        $4$ & $2$ \\
        \hline
        $5$ & $4$ \\
        \hline
        $6$ & $2$
    \end{tabular}
\end{table}
	
\section{Some section...}
\label{sec:some-section}

\lipsum
%\lipsum[3-56]
